%******************************************************************************
% preamble
%******************************************************************************
% type/format
\documentclass[12pt]{article}
\usepackage{geometry}
\geometry{
    letterpaper,
    margin=1in
}
% spacing | 1.6 === 2.0
\renewcommand{\baselinestretch}{1.6}
% font
\usepackage{fontspec}
\setmainfont{Times New Roman}
% language & encoding
\usepackage[spanish]{babel}
\usepackage[utf8]{inputenc}
% support images
\usepackage{graphicx}

% header & footer
\usepackage{fancyhdr}
\pagestyle{fancy}
\fancyhf{}
\rhead{\thepage}
\lfoot{\includegraphics[scale=0.40]{$HOME/downloads/espol-logo_footer.png}}

% title/author/date
\newcommand{\cTitle}{\LARGE{Formato IEEE-754}}
\newcommand{\cInstitution}{\textbf{Escuela Superior Politécnica del Litoral}}
\newcommand{\cClass}{Computación y Sociedad}
\newcommand{\cClassID}{CCPG1039}
\newcommand{\cTerm}{\textbf{I PAO - 2020}}
\newcommand{\cInstructor}{Lissette Cabello}
\title{
    \centering
    \includegraphics[scale=0.50]{$HOME/downloads/espol-logo_title.png}\\
    \textbf{\cTitle}
    \vspace{1in}
}
\author{
    \emph{Juan Antonio González Orbe}\\
    \normalsize \cInstitution\\
    \normalsize \textbf{\cClassID}: \cClass\\
    \normalsize \textbf{Instructor}: \cInstructor\\
    \normalsize \cTerm
}
\date{
    \vspace{2in}
    \centering
    \includegraphics[scale=0.15]{$HOME/downloads/espol_fiec-logo.png}\\
    \emph{\today}
}

% citation style
\usepackage[
backend=biber,
style=ieee,
citestyle=numeric
]{biblatex}
\addbibresource{$HOME/lib/docs/biblio.bib}

%******************************************************************************
% document itself
%******************************************************************************
\begin{document}
% title
\maketitle
% TOC
\clearpage
\tableofcontents
\clearpage

\section{Formato estándar IEEE-754}%
\label{sec:formato_estándar_ieee_754}

El formato \emph{IEEE-754} es un estándar técnico para representar números en
la notación \textit{coma flotante}. Esta norma fue constituida en el año 1985
por la \textit{IEEE}, que por sus siglas en inglés significa
``\textit{Institute of Electrical Electronics Engineers}''.
El propósito de este formato fue resolver distintos problemas que existían en
otras implementaciones de la notación \textit{coma flotante} que habían sido
desarrolladas en el pasado, entre algunos de los problemas estaba presente
la portabilidad. \cite{ieee754}

\subsection{Representación}%
\label{sub:representación}

El formato \emph{IEEE} se puede representa con $32-bits$ de los cuales, 
$23 bits$ forman la \textit{mantisa}, $8 bits$ forman el \textit{exponente} y 
el $1 bit$ que resta simboliza el \textit{signo} $+/-$ (positivo o negativo)
del número.

Sea un número ejemplar $11.87$ (positivo), el primer \textit{bit} será $0$,
mientras que si el número es negativo, sea $-45.36$, el primer \textit{bit}
será $1$. \cite{dec2ieee754yt}

\subsubsection{Ejemplo}%
\label{ssub:ejemplo}

Sea el número $-121.5$ para ejemplificar, se demostrará su valor en el formato
\textit{IEEE-74} que es $11000010111100110000000000000000$.

\begin{enumerate}
    \setlength{\parskip}{-8pt}
    \item Obtener el primer \textit{bit} 
        \begin{itemize}
            \setlength{\parskip}{-8pt}
            \item El primer \textit{bit} será $1$ ya que el número es negativo
        \end{itemize}
    \item Transformar el número a binario
        \begin{itemize}
            \setlength{\parskip}{-8pt}
            \item $-121.5$ en binario es $-1111001.1$
        \end{itemize}
    \item Correr la coma hasta tener solo $1$ número en la parte entera
        \begin{itemize}
            \setlength{\parskip}{-8pt}
            \item La coma se desplazó $6$ unidades
            \item $-1.1110011$
        \end{itemize}
    \item Transformar el exponente a binario
        \begin{itemize}
            \setlength{\parskip}{-8pt}
            \item Se suma $6 + 127$, que es el número de conversión para el
                formato \textit{IEEE-754}, resultando en $133$
            \item $133$ en binario es $10000101$
        \end{itemize}
    \item Obtener la \textit{mantisa} 
        \begin{itemize}
            \setlength{\parskip}{-8pt}
            \item Teniendo en cuenta el número obtenido anteriormente, 
                $-1.1110011$, se consideran los números a partir de la coma,
                ésto es $1110011$
            \item Se completa con $0$'s hasta completar $23$ dígitos.
            \item La \textit{mantisa} es $11100110000000000000000$
        \end{itemize}
    \item Obtener el número en el formato \textit{IEEE-754} 
        \begin{itemize}
            \setlength{\parskip}{-8pt}
            \item Se concatena el primer \textit{bit} más los \textit{8 bits} 
                del exponente más los \textit{23 bits} restantes de la mantisa
            \item Resultando en $11000010111100110000000000000000$
        \end{itemize}
\end{enumerate}

\section{Tarea}%
\label{sec:tarea}

Se solicita que cada estudiante realice la conversión al formato 
\textit{IEEE-754} a partir de 2 números decimales asignados por el instructor.

\subsection{Primer número}%
\label{sub:primer_número}

Sea $0.774$ el número solicitado:
\begin{enumerate}
    \setlength{\parskip}{-8pt}
    \item Obtener el primer \textit{bit} 
        \begin{itemize}
            \setlength{\parskip}{-8pt}
            \item El primer \textit{bit} será $0$ ya que el número es positivo
        \end{itemize}
    \item Transformar el número a binario
        \begin{itemize}
            \setlength{\parskip}{-8pt}
            \item $0.774$ en binario es $0.1100011000100100111$
        \end{itemize}
    \item Correr la coma hasta tener solo $1$ número en la parte entera
        \begin{itemize}
            \setlength{\parskip}{-8pt}
            \item La coma se desplazó $1$ unidad hacia la derecha ($-1$)
            \item $1.100011000100100111$
        \end{itemize}
    \item Transformar el exponente a binario
        \begin{itemize}
            \setlength{\parskip}{-8pt}
            \item Se suma $-1 + 127$, que es el número de conversión para el
                formato \textit{IEEE-754}, resultando en $126$
            \item $126$ en binario es $01111110$
        \end{itemize}
    \item Obtener la \textit{mantisa} 
        \begin{itemize}
            \setlength{\parskip}{-8pt}
            \item Teniendo en cuenta el número obtenido anteriormente, 
                $1.100011000100100111$, se consideran los números a partir de 
                la coma, ésto es $100011000100100111$
            \item Se completa con $0$'s hasta completar $23$ dígitos.
            \item La \textit{mantisa} es $10001100010010011100000$
        \end{itemize}
    \item Obtener el número en el formato \textit{IEEE-754} 
        \begin{itemize}
            \setlength{\parskip}{-8pt}
            \item Se concatena el primer \textit{bit} más los \textit{8 bits} 
                del exponente más los \textit{23 bits} restantes de la mantisa
            \item Resultando en $00111111010001100010010011100000$
        \end{itemize}
\end{enumerate}

\subsection{Segundo número}%
\label{sub:segundo}

Sea $-0.851$ el número solicitado:
\begin{enumerate}
    \setlength{\parskip}{-8pt}
    \item Obtener el primer \textit{bit} 
        \begin{itemize}
            \setlength{\parskip}{-8pt}
            \item El primer \textit{bit} será $1$ ya que el número es negativo
        \end{itemize}
    \item Transformar el número a binario
        \begin{itemize}
            \setlength{\parskip}{-8pt}
            \item $0.851$ en binario es $0.1101100111011011001$
        \end{itemize}
    \item Correr la coma hasta tener solo $1$ número en la parte entera
        \begin{itemize}
            \setlength{\parskip}{-8pt}
            \item La coma se desplazó $1$ unidad hacia la derecha ($-1$)
            \item $1.101100111011011001$
        \end{itemize}
    \item Transformar el exponente a binario
        \begin{itemize}
            \setlength{\parskip}{-8pt}
            \item Se suma $-1 + 127$, que es el número de conversión para el
                formato \textit{IEEE-754}, resultando en $126$
            \item $126$ en binario es $01111110$
        \end{itemize}
    \item Obtener la \textit{mantisa} 
        \begin{itemize}
            \setlength{\parskip}{-8pt}
            \item Teniendo en cuenta el número obtenido anteriormente, 
                $1.101100111011011001$, se consideran los números a partir de 
                la coma, ésto es $101100111011011001$
            \item Se completa con $0$'s hasta completar $23$ dígitos.
            \item La \textit{mantisa} es $10110011101101100100000$
        \end{itemize}
    \item Obtener el número en el formato \textit{IEEE-754} 
        \begin{itemize}
            \setlength{\parskip}{-8pt}
            \item Se concatena el primer \textit{bit} más los \textit{8 bits} 
                del exponente más los \textit{23 bits} restantes de la mantisa
            \item Resultando en $10111111010110011101101100100000$
        \end{itemize}
\end{enumerate}

\clearpage
\printbibliography
\end{document}
